\documentclass[conference]{IEEEtran}
\IEEEoverridecommandlockouts
% The preceding line is only needed to identify funding in the first footnote. If that is unneeded, please comment it out.
\usepackage{cite}
\usepackage{amsmath,amssymb,amsfonts}
\usepackage{algorithmic}
\usepackage{graphicx}
\usepackage{textcomp}
\usepackage{xcolor}
\def\BibTeX{{\rm B\kern-.05em{\sc i\kern-.025em b}\kern-.08em
    T\kern-.1667em\lower.7ex\hbox{E}\kern-.125emX}}
\begin{document}

\title{会议论文标题*\\
{\footnotesize \textsuperscript{*}注：子标题不会在 Xplore 中被捕获，不应使用}
\thanks{在此处确定适用的资助机构。如果没有，请删除此项。}
}

\author{\IEEEauthorblockN{第一作者 姓名}
\IEEEauthorblockA{\textit{所在部门} \\
\textit{机构名称}\\
城市, 国家 \\
电子邮件或 ORCID}
\and
\IEEEauthorblockN{第二作者 姓名}
\IEEEauthorblockA{\textit{所在部门} \\
\textit{机构名称}\\
城市, 国家 \\
电子邮件或 ORCID}
\and
\IEEEauthorblockN{第三作者 姓名}
\IEEEauthorblockA{\textit{所在部门} \\
\textit{机构名称}\\
城市, 国家 \\
电子邮件或 ORCID}
\and
\IEEEauthorblockN{第四作者 姓名}
\IEEEauthorblockA{\textit{所在部门} \\
\textit{机构名称}\\
城市, 国家 \\
电子邮件或 ORCID}
\and
\IEEEauthorblockN{第五作者 姓名}
\IEEEauthorblockA{\textit{所在部门} \\
\textit{机构名称}\\
城市, 国家 \\
电子邮件或 ORCID}
\and
\IEEEauthorblockN{第六作者 姓名}
\IEEEauthorblockA{\textit{所在部门} \\
\textit{机构名称}\\
城市, 国家 \\
电子邮件或 ORCID}
}

\maketitle

\begin{abstract}
\end{abstract}

\begin{IEEEkeywords}
\end{IEEEkeywords}

\section{引言}

Agent技能越来越多地用于为基于LLM的编码Agent打包可重用的能力。在实践中，一种技能很少只是一个单一的提示文件。它通常结合了指令文本、辅助脚本、配置、示例和外部资源。这使其看起来像一个小型软件工件，但其界面部分是自然语言，且其操作边界与Agent运行时行为纠缠在一起。本文的核心前提是，应该从软件组合角度和安全评估角度对技能进行研究。没有前者，安全讨论会过于抽象；没有后者，结构性观察将脱离实际的攻击和防御问题。

本文围绕三个研究问题展开。第一个问题探讨技能在野外是如何构成的，以及它们的关键组件是什么。这个问题旨在建立一个稳定的描述性基准。第二个问题探讨与技能相关的攻击向量和攻击面。这个问题旨在通过将先前关于提示注入、供应链滥用和Agent工具的工作与技能工件的具体结构联系起来，来细化当前对技能安全的理解。第三个问题探讨当前攻防之间的博弈现状。这个问题从描述和分类转向对比评估，利用代表性攻击和现有扫描器来理解技能工件的哪些部分得到了很好的覆盖，哪些部分仍然防御薄弱。

因此，全文将分三层展开。它将首先表征该工件。然后，它将围绕该工件组织威胁空间。最后，它将评估当针对不同技能组件实例化攻击时，现有防御工具的表现。在当前阶段，本草案的目的是为这三层确定一个清晰且可编写的大纲，以便后续迭代可以填充测量数据、表格和讨论，而无需反复更改整体论点。

\section{背景与范围}

本节将技能定义为一种独特的软件工件。第一部分将定义实践中什么构成了技能，包括指令文件（如 \texttt{SKILL.md}）、辅助脚本、元数据、清单（manifests）、嵌入式示例以及对远程资源的引用。目标是表明技能既不是纯粹的提示，也不是纯粹的包。它是一个混合工件，其组合方式对可用性和安全性都至关重要。

本节的第二部分将明确本文的研究范围。焦点将集中在由技能工件本身携带的攻击，或者由工件加载和使用方式直接触发的攻击。这包括恶意指令、欺骗性文档、脚本级有效载荷、误导性声明，以及利用工件内容与Agent操作之间关系的攻击。相比之下，本文不会涉及模型预训练攻击、与技能无关的通用Web提示注入，或不依赖于技能渠道的广泛运行时安全问题。这种界定很重要，因为Agent安全周围的空间已经很大，本文需要一个清晰的重点。

\section{方法论}

方法论部分将解释三个研究问题如何融入到一个连贯的研究设计中。研究始于语料库构建。我们计划从多个渠道收集大约一千个技能，包括公开市场和基于仓库的发行版，同时保持流行度和发布者类型的多样性。目标不是建立生态系统的完整普查，而是构建一个足够广泛的语料库，以揭示反复出现的结构模式。

在语料库构建之后，本文将对采样的技能进行结构分析。这部分支持 RQ1。我们将检查文件类型、目录布局、脚本的存在、配置工件的存在，以及对外部依赖或服务的引用。然后，我们将这些观察结果与传统的软件架构概念（如接口、实现、依赖和打包）进行对比。这一步很重要，因为本文不仅在问技能包含什么，还在问应该如何解释这种构成。

下一个方法论组件是轻量级的文献综述，用以支持 RQ2。综述将以近期关于技能安全、Agent安全、提示注入和可转移供应链攻击的研究作为主要搜索领域，随后对高度相关的论文和工具进行一轮向后滚雪球式搜索。目的不是进行完整的系统综述，而是建立一个可辩护的、分层的攻击向量说明，该说明需充分基于先前工作，同时与 RQ1 中观察到的技能工件直接挂钩。

最后一个方法论组件是 RQ3 的评估协议。现有的基准测试和实验结果将作为起点。我们将根据攻击目标组件对现有攻击案例进行整理，并在必要时添加操纵变体，将相同的恶意意图跨指令、脚本、元数据和外部引用进行移动。现有的扫描器和模型驱动的分析器将在统一的框架下运行。主要产出将是当前防御措施在当前设置下能检测和不能检测什么的对比视图，以及对这些行为在多大程度上可以由技能组合本身解释的讨论。

\section{RQ1: 技能作为软件工件是如何构成的}

第一个研究问题旨在为本文提供描述性基础。本节将首先量化采样语料库中出现了哪些主要的组件。它将检查技能仅包含文档的频率、包含可执行文件的频率、附带支持示例或引用的频率，以及声明或隐含与外部工具和服务交互的频率。本节不会将这些观察结果呈现为杂乱无章的清单，而是利用它们来识别少数反复出现的结构原型。

本节的第二部分将从软件工件的角度解读这些原型。文档密集的技能可能更像是一个自然语言接口封装器。具有多个脚本和辅助文件的技能可能类似于轻量级插件包。对外部端点有强依赖的技能可能表现得更像是一个配置存根（stub），将关键行为委托给其他地方。该分析的重点是解释在技能背景下，哪些部分保留了传统软件打包的特征，哪些部分是真正创新的。

RQ1 的最后部分将通过论证工件的组成不仅是描述性的背景，来为后续章节做铺垫。它也是攻击如何隐藏、权限如何被误报以及扫描器可能失效的基础。从这个意义上说，RQ1 是本文其余部分赖以生存的结构层。

\section{RQ2: 技能的攻击向量和攻击面是什么}

第二个研究问题将把当前对技能攻击的扁平化描述转化为更明确的层次结构。本节将首先回顾最相关的先前工作，并将攻击知识分为三个来源：专门为技能设计的攻击；最初为智能体或提示注入设计但可直接转移到技能的攻击；以及改编自传统软件供应链的攻击。这种区分是必要的，因为本文不应将所有攻击向量都视为技能设置中所特有的。

本节的下一部分将通过将攻击目标与工件组件及触发条件联系起来，来组织攻击空间。在宏观层面上，分类学将区分攻击嵌入的位置、它试图获取的能力以及它在什么条件下生效。这种层次结构旨在改进简单的攻击名称列表。它不仅应展示存在多个攻击家族，还应展示它们如何与 RQ1 中识别的组合模式相关联。

RQ2 的最后部分将解释该分类学将如何在该论文后续部分中使用。它将作为语料库表征与实证评估之间的概念桥梁。它将定义哪些类型的攻击放置测试是有意义的，可以合理地提出哪些类型的扫描器覆盖声明，以及当前文献在哪些方面仍显不足。这使得 RQ2 不仅仅是一个综述章节，它是后续评估的组织理论。

\section{RQ3: 当前的攻防态势如何}

第三个研究问题是本文的实证对抗部分。它将选取具有代表性的攻击，并将其置于现有的扫描器和模型驱动分析器面前。目的不仅是产生一个排行榜，而是观察当攻击放置在技能工件的不同部分时，防御行为会如何变化，以及某些技能组件是否始终比其他组件更容易或更难防御。

本节将首先描述已有的攻击集以及如何组织它们进行评估。现有的攻击案例提供了一个切实的起点，但本节也将为操纵变体留出空间，这些变体在保持相同恶意意图的同时改变载体组件。例如，最初在指令文本中表达的攻击可能会通过辅助脚本或误导性元数据重新呈现。这种设计将允许本文将实证结果联系回 RQ1 中开发的工件模型和 RQ2 中开发的攻击层次结构。

下一部分将描述评估目标。计划在通用协议下测试多个扫描器和模型驱动的设置，包括在可行的情况下测试能力较低的基准线。核心问题不仅是哪个系统能检测到更多的案例，而且是在相同的意图被移动到不同组件后，哪些类型的攻击仍然可以被检测到。即使绝对扫描器性能证明比最初预期的更强，本文也可以在此提出有用的主张。

RQ3 的最后部分将把结果界定为一场“博弈”而非静态的基准报告。它将讨论现有防御在哪些方面表现稳健，在哪些方面对组件敏感，以及当前的基准在哪些方面仍然太弱而无法支持强有力的主张。因此，本节应以平衡的解释结束：要么证据支持一篇有意义的安全评估论文，要么显现出该贡献较窄，应更审慎地界定其范围。

\section{讨论}

\section{结论}

\section*{致谢}


\section*{参考文献}

\begin{thebibliography}{00}
\bibitem{placeholder} 参考文献将在实证部分和文献综述定稿后插入。
\end{thebibliography}

\end{document}
